\chapter{Sprint 4 : Service d'inférence et poste de travail}
\label{chap:sprint4}

\sectionsansnum{Introduction}

Les trois premiers sprints ont produit des modèles. Ce quatrième sprint les rend
utilisables : il expose une interface applicative sécurisée et construit le poste de
travail du praticien, depuis le dépôt d'un enregistrement jusqu'au rapport de sévérité
expliqué.

\section{Objectif du sprint}

L'objectif était de franchir le pas du \emph{carnet d'expérimentation} vers le
\emph{service}. Trois exigences le structuraient.

La première était la \textbf{disponibilité des modèles}. Chargés naïvement à chaque
requête, les modèles imposeraient plusieurs secondes de latence : ils devaient donc être
conservés en mémoire entre deux appels.

La deuxième était la \textbf{sécurisation des accès}. Une plateforme manipulant des
données de santé ne peut être ouverte : authentification, jetons signés et cloisonnement
des rôles devaient être posés dès ce sprint.

La troisième était l'\textbf{ergonomie du parcours}. Une analyse dure plusieurs dizaines
de secondes ; l'interface devait guider le praticien et le tenir informé sans jamais le
laisser devant un écran muet.

\section{Analyse du sprint 4}

\subsection{Sprint Backlog}

Le sprint a embarqué seize récits pour 80 points, le plus chargé du projet.

\begin{table}[H]
\centering\small
\renewcommand{\arraystretch}{1.15}
\begin{tabular}{@{}c|L{2.5cm}|L{4.6cm}|L{4.6cm}|c@{}}
\hline
\ligneentete \entete{ID} & \entete{Module} & \entete{R\'ecit utilisateur} & \entete{T\^aches} & \entete{Estim.} \\
\hline
\multirow{3}{*}{US-01\,\`a\,04} & \multirow{3}{=}{Authentification} & \multirow{3}{=}{En tant qu'utilisateur, je veux me connecter et acc\'eder aux fonctions de mon r\^ole} & Hacher les mots de passe avec sel al\'eatoire & \multirow{3}{*}{15 pts} \\
 &  &  & D\'elivrer un jeton sign\'e \`a expiration &  \\
 &  &  & V\'erifier le r\^ole sur chaque route prot\'eg\'ee &  \\
\hline
\multirow{3}{*}{US-30\,\`a\,33} & \multirow{3}{=}{Service d'inf\'erence} & \multirow{3}{=}{En tant que m\'edecin, je veux lancer une analyse en choisissant la configuration adapt\'ee \`a mon montage} & Exposer les routes d'analyse et d'inspection & \multirow{3}{*}{19 pts} \\
 &  &  & Mettre les mod\`eles en cache m\'emoire &  \\
 &  &  & Calculer les m\'etriques AASM &  \\
\hline
\multirow{3}{*}{US-38\,\`a\,41} & \multirow{3}{=}{Rapport de s\'ev\'erit\'e} & \multirow{3}{=}{En tant que m\'edecin, je veux saisir les donn\'ees cliniques et obtenir un rapport expliqu\'e} & Cr\'eer le formulaire clinique & \multirow{3}{*}{23 pts} \\
 &  &  & Analyser les rapports externes import\'es &  \\
 &  &  & Produire l'interpr\'etation automatique &  \\
\hline
\multirow{3}{*}{US-42\,\`a\,44} & \multirow{3}{=}{Visualisation} & \multirow{3}{=}{En tant que m\'edecin, je veux visualiser l'hypnogramme et suivre la progression de l'analyse} & Tracer l'hypnogramme interactif & \multirow{3}{*}{18 pts} \\
 &  &  & Superposer saturation et densit\'e d'\'ev\'enements &  \\
 &  &  & Afficher l'\'ecran de progression &  \\
\hline
\multirow{1}{*}{US-53} & \multirow{1}{=}{Export} & \multirow{1}{=}{En tant que m\'edecin, je veux exporter les marqueurs calcul\'es} & Produire les exports tabulaire et balis\'e & \multirow{1}{*}{5 pts} \\
\hline
\end{tabular}
\caption{Sprint Backlog --- Sprint 4}
\label{tab:bl-s4}
\end{table}

\subsection{Diagramme des cas d'utilisation}

\begin{figure}[H]
\centering
\begin{tikzpicture}[
    x=1cm,y=1cm,
    uc/.style     = {ellipse, draw=hypnoblue!60, fill=hypnolight,
                     align=center, font=\scriptsize, text width=2.7cm,
                     inner sep=1.5pt, minimum height=1.0cm},
    ucp/.style    = {ellipse, draw=hypnored!70, fill=hypnored!7,
                     align=center, font=\scriptsize\bfseries, text width=2.7cm,
                     inner sep=1.5pt, minimum height=1.0cm},
    lien/.style   = {hypnogray!75, line width=0.55pt},
    rel/.style    = {-{Latex[length=2mm]}, hypnoblue!70, line width=0.5pt, dashed},
    et/.style     = {font=\tiny\itshape, text=hypnoblue!85, fill=white, inner sep=1pt},
    nomact/.style = {font=\scriptsize\bfseries, align=center, text width=2.2cm},
  ]

  \tikzset{acteur/.pic={
      \draw[hypnoblue,line width=0.8pt] (0,0) circle (0.17);
      \draw[hypnoblue,line width=0.8pt] (0,-0.17) -- (0,-0.66);
      \draw[hypnoblue,line width=0.8pt] (-0.28,-0.34) -- (0.28,-0.34);
      \draw[hypnoblue,line width=0.8pt] (0,-0.66) -- (-0.24,-1.04);
      \draw[hypnoblue,line width=0.8pt] (0,-0.66) -- (0.24,-1.04);}}

  \draw[hypnoblue!55,line width=0.9pt,rounded corners=3pt]
       (3.6,-5.4) rectangle (12.6,2.4);
  \node[font=\small\bfseries,text=hypnoblue,anchor=north west] at (3.85,2.25)
       {Poste de travail du médecin};

  \node[uc]  (auth) at (5.8, 1.4)  {S'authentifier};
  \node[ucp] (ana)  at (5.8,-0.6)  {Analyser un\\enregistrement};
  \node[uc]  (exp)  at (5.8,-3.9)  {Exporter les\\marqueurs};

  \node[uc]  (cfg)  at (10.3, 1.3) {Choisir montage\\et granularité};
  \node[uc]  (vis)  at (10.3,-0.2) {Visualiser\\l'hypnogramme};
  \node[uc]  (met)  at (10.3,-1.7) {Consulter les\\métriques AASM};
  \node[uc]  (sev)  at (10.3,-3.2) {Obtenir le rapport\\de sévérité};
  \node[uc]  (imp)  at (10.3,-4.7) {Importer un\\rapport externe};

  \pic at (1.5,0.9) {acteur};
  \node[nomact,anchor=north] at (1.5,0.25) {Médecin};

  \draw[lien] (1.9,0.55) -- (auth.west);
  \draw[lien] (1.9,0.55) -- (ana.west);
  \draw[lien] (1.9,0.55) -- (exp.west);

  \draw[rel] (ana) -- node[et,pos=0.58] {$\ll$include$\gg$} (cfg);
  \draw[rel] (ana) -- node[et,pos=0.6]  {$\ll$include$\gg$} (vis);
  \draw[rel] (ana) -- node[et,pos=0.6]  {$\ll$include$\gg$} (met);
  \draw[rel] (sev) -- node[et,pos=0.5,above]  {$\ll$extend$\gg$} (met);
  \draw[rel] (imp) -- node[et,pos=0.5,below]  {$\ll$extend$\gg$} (sev);
  \draw[rel] (ana) -- node[et,pos=0.55] {$\ll$include$\gg$} (auth);

\end{tikzpicture}
\caption{Diagramme des cas d'utilisation du sprint 4 --- poste de travail du médecin}
\label{fig:uc-s4}
\end{figure}

\subsubsection{Description textuelle des cas d'utilisation}

\begin{casutilisation}{Description textuelle du cas « Analyser un enregistrement »}
Titre & Analyser un enregistrement \\ \hline
Acteurs & Médecin du sommeil, moteur d'inférence \\ \hline
Description & Produire l'hypnogramme et les métriques d'une nuit à partir d'un fichier
d'enregistrement brut \\ \hline
Préconditions & Le médecin est authentifié ; sa licence est valide ; il dispose d'un
fichier au format EDF \\ \hline
Scénario de base &
\begin{minipage}[t]{\linewidth}
\begin{enumerate}[leftmargin=1.1em,itemsep=1pt,topsep=3pt]
  \item Le médecin choisit le montage disponible, à deux ou cinq dérivations.
  \item Il choisit la granularité de classification, à trois ou cinq stades.
  \item Il dépose le fichier d'enregistrement.
  \item Le système décode le fichier, résout les dérivations et prétraite le signal.
  \item Le moteur classe chaque époque et le système calcule les métriques.
  \item Le système affiche l'hypnogramme, les métriques et les alertes cliniques.
\end{enumerate}
\end{minipage} \\ \hline
Scénario alternatif &
\begin{minipage}[t]{\linewidth}
\begin{itemize}[leftmargin=1.1em,itemsep=1pt,topsep=3pt]
  \item \textbf{3a} --- Le fichier n'est pas au format attendu : le dépôt est refusé.
  \item \textbf{4a} --- Une dérivation est absente : repli documenté, l'analyse aboutit.
  \item \textbf{4b} --- L'enregistrement est trop court : interruption avec message.
\end{itemize}
\end{minipage}
\end{casutilisation}

\subsection{Diagramme de séquence objet}

\begin{figure}[H]
\centering
\resizebox{0.98\textwidth}{!}{%
\begin{sequencediagram}
  \newthread[hypnolight]{u}{\scriptsize Praticien}
  \newinst[2.2]{f}{\scriptsize Client web}
  \newinst[2.2]{a}{\scriptsize Interface applicative}
  \newinst[2.2]{d}{\scriptsize Base de données}

  \begin{call}{u}{\scriptsize saisit ses identifiants}{f}{}
  \end{call}
  \begin{call}{f}{\scriptsize demande de jeton}{a}{\scriptsize jeton et profil}
    \begin{call}{a}{\scriptsize recherche du compte}{d}{\scriptsize empreinte}
    \end{call}
    \begin{sdblock}{\scriptsize Vérifications}{\scriptsize mot de passe, statut, licence}
      \begin{call}{a}{\scriptsize horodate la connexion}{d}{}
      \end{call}
      \begin{call}{a}{\scriptsize journalise l'accès}{d}{}
      \end{call}
    \end{sdblock}
  \end{call}
  \begin{call}{f}{\scriptsize conserve le jeton}{f}{}
  \end{call}
  \begin{call}{f}{\scriptsize ouvre l'espace de travail}{u}{}
  \end{call}
\end{sequencediagram}}
\caption[Séquence d'authentification]%
        {Ouverture de session et délivrance du jeton.}
\label{fig:seq-auth}
\end{figure}

\begin{figure}[H]
\centering
\resizebox{0.98\textwidth}{!}{%
\begin{sequencediagram}
  \newthread[hypnolight]{u}{\scriptsize Praticien}
  \newinst[1.8]{f}{\scriptsize Client web}
  \newinst[1.8]{a}{\scriptsize Interface applicative}
  \newinst[1.8]{m}{\scriptsize Moteur d'inférence}
  \newinst[1.8]{d}{\scriptsize Base de données}
  \newinst[1.8]{s}{\scriptsize Stockage objet}

  \begin{call}{u}{\scriptsize dépose le fichier}{f}{}
  \end{call}
  \begin{call}{f}{\scriptsize demande d'analyse}{a}{\scriptsize stades et métriques}
    \begin{call}{a}{\scriptsize prétraitement}{m}{\scriptsize tenseur d'époques}
    \end{call}
    \begin{call}{a}{\scriptsize classification}{m}{\scriptsize séquence de stades}
    \end{call}
    \begin{call}{a}{\scriptsize métriques AASM}{m}{\scriptsize indicateurs}
    \end{call}
  \end{call}
  \begin{call}{f}{\scriptsize affiche l'hypnogramme}{u}{}
  \end{call}
  \begin{call}{f}{\scriptsize enregistre l'examen}{a}{\scriptsize identifiant}
    \begin{call}{a}{\scriptsize insertion}{d}{}
    \end{call}
  \end{call}
  \begin{sdblock}{\scriptsize Arrière-plan}{\scriptsize sans blocage de l'interface}
    \begin{call}{f}{\scriptsize transfert du fichier}{a}{\scriptsize référence}
      \begin{call}{a}{\scriptsize dépôt}{s}{\scriptsize clé d'objet}
      \end{call}
      \begin{call}{a}{\scriptsize mise à jour}{d}{}
      \end{call}
    \end{call}
  \end{sdblock}
\end{sequencediagram}}
\caption[Séquence d'analyse d'un enregistrement]%
        {Analyse d'un enregistrement : les résultats précèdent l'archivage.}
\label{fig:seq-analyse}
\end{figure}

La particularité de cet enchaînement tient à l'ordre des opérations : les résultats sont
affichés au praticien \emph{avant} que l'archivage du fichier, mis en place au sprint
suivant, ne soit terminé.

\section{Réalisation}

\subsection{Organisation de l'interface applicative}

L'interface applicative est bâtie sur FastAPI, choisi pour son modèle asynchrone, sa
validation déclarative des données et la génération automatique de sa documentation. Le
code est réparti en modules à responsabilité unique.

\begin{table}[H]
\centering

\small
\begin{tabular}{@{}lrL{8.3cm}@{}}
\toprule
\textbf{Module} & \textbf{Lignes} & \textbf{Responsabilité} \\
\midrule
\texttt{main.py}          & 695 & Point d'entrée, routes métier et d'administration,
                                  migration au démarrage, journalisation \\
\texttt{ml\_routes.py}    & 984 & Routes d'inférence : analyse, extraction de
                                  marqueurs, prédiction de sévérité, lecture de
                                  fichiers externes \\
\texttt{ml\_models.py}    & 117 & Définition des trois architectures de réseaux \\
\texttt{preprocessing.py} & 199 & Décodage EDF, résolution des dérivations, calcul
                                  des métriques AASM \\
\texttt{osa\_predictor.py}& 113 & Chargement des modèles de sévérité, prédiction et
                                  attribution des contributions \\
\texttt{db\_models.py}    & 113 & Modèle de données et relations \\
\texttt{schemas.py}       & 163 & Schémas de validation des entrées et sorties \\
\texttt{auth.py}          &  51 & Jetons, empreintes de mots de passe, dépendance
                                  d'authentification \\
\texttt{b2\_storage.py}   & 121 & Client de stockage objet, dépôt et signature des
                                  liens \\
\texttt{database.py}      &  19 & Connexion et session \\
\bottomrule
\end{tabular}
\caption{Modules de l'interface applicative}
\label{tab:modules}
\end{table}

Deux mécanismes transversaux méritent d'être signalés. Le premier est la \textbf{migration
au démarrage} : le schéma ayant évolué au fil des sprints, l'application exécute à son
lancement une série d'ajouts de colonnes conditionnels, ce qui permet de déployer une
nouvelle version sur une base existante sans intervention manuelle. Le second est la
\textbf{signature des liens à la sérialisation}, développée au sprint suivant.

\subsection{Points d'entrée}

\begin{table}[H]
\centering

\small
\begin{tabular}{@{}lL{6.6cm}cL{2.6cm}@{}}
\toprule
\textbf{Méthode} & \textbf{Chemin} & \textbf{Accès} & \textbf{Domaine} \\
\midrule
POST & \texttt{/token}                        & public   & \multirow{2}{*}{Authentification} \\
GET  & \texttt{/users/me}                     & connecté & \\
\midrule
POST & \texttt{/admin/doctors}                & admin & \multirow{8}{*}{Administration} \\
GET  & \texttt{/admin/doctors}                & admin & \\
PUT  & \texttt{/admin/doctors/\{id\}/lifecycle}      & admin & \\
POST & \texttt{/admin/doctors/\{id\}/reset-password} & admin & \\
POST & \texttt{/admin/hospitals}              & admin & \\
GET  & \texttt{/admin/hospitals}              & admin & \\
GET  & \texttt{/admin/dashboard-stats}        & admin & \\
GET  & \texttt{/admin/audit-logs}             & admin & \\
GET  & \texttt{/admin/audit-logs/export}      & admin & \\
\midrule
POST & \texttt{/patients}                     & médecin & \multirow{4}{*}{Patients} \\
GET  & \texttt{/patients}                     & médecin & \\
GET  & \texttt{/patients/\{id\}}              & médecin & \\
GET  & \texttt{/doctor/stats}                 & médecin & \\
\midrule
POST & \texttt{/patients/\{id\}/psgs}         & médecin & \multirow{8}{*}{Examens} \\
PUT  & \texttt{/psgs/\{id\}}                  & médecin & \\
POST & \texttt{/psgs/\{id\}/upload\_edf}      & médecin & \\
POST & \texttt{/psgs/\{id\}/upload\_hypnogram} & médecin & \\
POST & \texttt{/psgs/\{id\}/upload\_hypnogram\_annotated} & médecin & \\
POST & \texttt{/psgs/\{id\}/upload\_osa\_report} & médecin & \\
POST & \texttt{/psgs/\{id\}/annotations}      & médecin & \\
GET  & \texttt{/psgs/\{id\}/annotations}      & médecin & \\
\midrule
GET  & \texttt{/doctors}                      & médecin & \multirow{5}{*}{Collaboration} \\
POST & \texttt{/conversations}                & médecin & \\
GET  & \texttt{/conversations}                & médecin & \\
GET  & \texttt{/conversations/psg/\{id\}}     & médecin & \\
GET  & \texttt{/conversations/\{id\}/messages} & médecin & \\
POST & \texttt{/conversations/\{id\}/messages} & médecin & \\
\midrule
GET  & \texttt{/health}                       & public & \multirow{7}{*}{Inférence} \\
POST & \texttt{/channels}                     & public & \\
POST & \texttt{/analyze}                      & public & \\
POST & \texttt{/extract\_features}            & public & \\
POST & \texttt{/predict\_osa}                 & public & \\
POST & \texttt{/predict\_osa\_custom}         & public & \\
POST & \texttt{/parse\_features\_file}        & public & \\
\bottomrule
\end{tabular}
\caption{Points d'entrée de l'interface applicative}
\label{tab:endpoints}
\end{table}

Il faut relever ici une \textbf{faiblesse de la version livrée} : les routes d'inférence
ne sont pas protégées par le contrôle d'accès. Ce choix, initialement fait pour faciliter
les essais pendant la mise au point, n'a pas été rétabli. Un appelant non authentifié peut
donc soumettre un enregistrement à l'analyse. Aucune donnée patient n'est exposée --- ces
routes ne lisent ni n'écrivent en base --- mais elles consomment des ressources de calcul.
L'ajout de la dépendance d'authentification, qui tient en une ligne par route, figure parmi
les corrections à apporter avant toute mise en service réelle.

\capturefig{swagger}{0.9}{Documentation OpenAPI générée automatiquement}{fig:swagger}

\subsection{Sécurité des accès}

\begin{table}[H]
\centering

\small
\begin{tabular}{@{}L{3.2cm}L{5.3cm}L{5.5cm}@{}}
\toprule
\textbf{Mécanisme} & \textbf{Principe} & \textbf{Menace couverte} \\
\midrule
Empreintes de mots de passe &
Fonction de dérivation lente avec sel aléatoire par compte &
Compromission de la base : les mots de passe restent inexploitables \\
\addlinespace[3pt]
Jetons signés &
Jeton porteur signé, transmis à chaque appel, contenant l'identité et le rôle &
Usurpation : toute altération invalide la signature \\
\addlinespace[3pt]
Contrôle d'accès par rôle &
Vérification du rôle et de la propriété de la ressource sur chaque route protégée &
Élévation de privilège et accès aux patients d'un confrère \\
\addlinespace[3pt]
Liens d'accès signés &
Les références de fichiers sont converties en URL signées valables 24 heures &
Fuite de données : aucun objet n'est accessible sans signature valide \\
\bottomrule
\end{tabular}
\caption{Mécanismes de sécurité et menaces couvertes}
\label{tab:secu}
\end{table}

Les \textbf{mots de passe} sont hachés par bcrypt, dont le sel aléatoire par compte
interdit les tables précalculées. Les \textbf{jetons} sont signés en HMAC-SHA256 et portent
l'adresse de l'utilisateur, son rôle et une date d'expiration fixée à sept jours. Le
\textbf{contrôle d'accès} superpose deux niveaux : le rôle est vérifié en tête de route, la
propriété de la ressource l'est ensuite.

\subsection{Poste de travail du praticien}

L'interface est une application monopage React, construite par Vite, organisée en
composants regroupés par domaine fonctionnel.

\begin{table}[H]
\centering

\small
\begin{tabular}{@{}lrL{8.0cm}@{}}
\toprule
\textbf{Composant} & \textbf{Lignes} & \textbf{Rôle} \\
\midrule
\texttt{AnalysisResults} & 1134 & Restitution complète : hypnogramme, métriques,
                                  marqueurs, formulaire clinique, rapport de sévérité \\
\texttt{PatientList}     &  993 & Registre des patients, historique des examens,
                                  documents archivés \\
\texttt{CustomOSA}       &  658 & Prédiction à partir d'une saisie manuelle ou d'un
                                  rapport importé \\
\texttt{Hypnogram}       &  552 & Tracé interactif de l'hypnogramme, de la saturation
                                  et de la densité d'événements \\
\texttt{DoctorList}      &  539 & Administration des comptes et des licences \\
\texttt{DeveloperPipeline} & 497 & Visualisation de la chaîne de traitement \\
\texttt{LandingPage}     &  406 & Page publique et authentification \\
\texttt{AnalysisWizard}  &  428 & Assistant d'analyse en quatre étapes \\
\texttt{HospitalsDashboard} & 233 & Gestion des établissements \\
\texttt{ConsultationsView}  & 198 & Boîte de réception des sollicitations \\
\texttt{AuditLogsDashboard} & 194 & Consultation et export du journal \\
\texttt{CollaborationChat}  & 171 & Fil de discussion rattaché à un examen \\
\bottomrule
\end{tabular}
\caption{Principaux composants de l'interface}
\label{tab:composants}
\end{table}

Le parcours principal se déroule en quatre étapes : le praticien déclare le montage
disponible, puis la granularité souhaitée, dépose le fichier, et consulte les résultats.
Chaque étape n'expose qu'une seule décision, assortie d'une recommandation par défaut.

\capturefig{wizard-1}{0.86}{Assistant d'analyse, étape 1 : déclaration du montage}{fig:wizard1}

\capturefig{wizard-2}{0.86}{Assistant d'analyse, étape 2 : granularité de classification}{fig:wizard2}

Pendant le traitement, un écran de progression affiche les tracés physiologiques et
détaille les cinq étapes du pipeline. Ce retour visuel n'est pas décoratif : une interface
muette pendant plusieurs dizaines de secondes laisserait croire à un blocage.

\capturefig{progression}{0.86}{Écran de progression du traitement}{fig:progression}

La page de résultats superpose sur un même axe temporel l'hypnogramme prédit, la courbe de
saturation et la densité d'événements respiratoires. Cette superposition est le choix de
visualisation le plus utile de l'interface : elle permet de constater directement qu'une
désaturation coïncide avec une sortie de sommeil paradoxal, ce que trois graphiques séparés
ne montreraient pas.

\capturefig{resultats}{0.92}{Restitution des résultats d'analyse}{fig:resultats}

Un second parcours permet une prédiction sans enregistrement, à partir d'une saisie manuelle
ou d'un rapport externe importé. Cette voie répond à un usage réel : disposer d'un compte
rendu produit par un autre laboratoire sans avoir accès au signal brut.

\capturefig{custom-osa}{0.86}{Prédiction de sévérité à partir d'une saisie manuelle}{fig:customosa}

\sectionsansnum{Conclusion}

Ce sprint a transformé des modèles de laboratoire en un service utilisable. L'interface
applicative expose les routes d'inférence, conserve les modèles en mémoire et protège les
routes métier par un double contrôle de rôle et de propriété. Le poste de travail guide
l'analyse en quatre décisions et restitue les résultats sur un axe temporel unique.

Une faiblesse a été identifiée et documentée : les routes d'inférence ne sont pas
authentifiées.

Le sprint suivant dote la plateforme d'une mémoire : dossier patient, archivage des fichiers
volumineux et administration multi-établissements.
